\chapter{Planificación y metodología de desarrollo}
\label{cap:planificacion}

Este capítulo describe cómo se organizó el proyecto. Presenta el enfoque iterativo, las fases, la dedicación, el calendario y los riesgos principales. La metodología experimental se desarrolla en el capítulo 8.

\section{Enfoque de trabajo}
\label{sec:enfoque_trabajo}

El desarrollo siguió un enfoque iterativo e incremental. Se partió de un prototipo limitado a \textit{IRS vanilla} y cada ciclo produjo una versión ejecutable que permitió revisar requisitos, esquema, instrucciones, referencias y validación.

Ante un resultado incorrecto se comprobaban tres cuestiones antes de atribuirlo al modelo: que la instrucción fuera clara, que el agente hubiera recibido la información necesaria y que la referencia fuera alcanzable. Esta práctica resultó esencial porque varios fallos iniciales procedían del instrumento de evaluación y no del modelo.

El código, las instrucciones, el esquema y los casos se gestionaron mediante control de versiones. Las ejecuciones se separaron en tandas identificables para no combinar resultados obtenidos bajo condiciones distintas. La configuración de modelos y agentes permaneció fuera de la lógica principal, lo que permitió comparar alternativas sin modificar las reglas financieras.

Este enfoque fue adecuado porque el problema reunía dos fuentes de incertidumbre: la definición financiera evolucionó al contrastar convenciones y el comportamiento de los modelos solo podía conocerse mediante ejecuciones. Los ciclos permitieron corregir supuestos sin perder la trazabilidad de las versiones anteriores.

\section{Fases del proyecto}
\label{sec:fases_proyecto}

El trabajo se dividió en siete fases relacionadas. No fueron completamente secuenciales: las pruebas y la evaluación llevaron a revisar decisiones anteriores. La Tabla~\ref{tab:fases_proyecto} resume las actividades y entregables.

\begin{table}[H]
\centering
\small
\caption{Fases principales del proyecto}
\label{tab:fases_proyecto}
\begin{tabular}{|P{3.0cm}|P{6.8cm}|P{3.8cm}|}
\hline
\textbf{Fase} & \textbf{Actividades} & \textbf{Resultado} \\
\hline
Análisis y alcance & Estudio del flujo, dificultades y límites del producto. & Problema y requisitos iniciales. \\
\hline
Representación financiera & Campos, patas, propósitos, curvas, convenciones y calendarios. & Modelo de la \textit{RFQ}. \\
\hline
Diseño & Responsabilidades, validación y contrato de datos. & Arquitectura de la solución. \\
\hline
Implementación & Agentes, modelos internos, validación, mapeo y proveedores. & Prototipo ejecutable. \\
\hline
Pruebas y refinamiento & Pruebas automatizadas, casos y revisión del instrumento. & Versión corregida y trazable. \\
\hline
Evaluación & Seis modelos sobre 23 casos, con telemetría y métricas. & Evidencia cuantitativa. \\
\hline
Documentación & Consolidación de decisiones, resultados y limitaciones. & Memoria y material de defensa. \\
\hline
\end{tabular}
\floatfoot{Fuente: elaboración propia.}
\end{table}

La fase financiera se apoyó en la estructura del \textit{IRS} descrita por \textcite[Definición~2.19, p.~59]{ausin2025quantitative} y en fuentes sectoriales para las convenciones. El diseño separó clasificación, extracción, validación y mapeo. Durante las pruebas se construyeron los casos y se revisaron las referencias antes de cerrar la comparación de modelos.

La evaluación produjo la evidencia final, pero también confirmó el valor del proceso iterativo: una instrucción ambigua, una entrada incompleta al agente o una referencia asimétrica podían parecer fallos del modelo. Las incidencias y sus resultados se analizan en el capítulo 9, no se repiten aquí.

\subsection{Recursos y dedicación por fase}
\label{sec:recursos_dedicacion}

El TFM corresponde a seis créditos ECTS y una dedicación de 180 horas. Las cinco horas de seguimiento presencial con el tutor se consideran integradas en las fases con revisiones financieras, metodológicas y documentales; no se añaden como una partida separada.

\begin{table}[H]
\centering
\small
\caption{Distribución estimada de dedicación y recursos}
\label{tab:dedicacion_fases}
\begin{tabular}{|P{3.1cm}|P{1.3cm}|P{5.0cm}|P{4.0cm}|}
\hline
\textbf{Fase} & \textbf{Horas} & \textbf{Recursos principales} & \textbf{Justificación} \\
\hline
Análisis y delimitación & 18 & Estudiante, tutor y bibliografía. & Definir problema, alcance y requisitos. \\
\hline
Representación financiera & 24 & Tutor, libro de referencia y fuentes sectoriales. & Diseñar el \textit{IRS} y sus convenciones. \\
\hline
Diseño de la solución & 20 & Python, protobuf y documentación técnica. & Definir arquitectura y validación. \\
\hline
Implementación & 32 & Equipo Windows, Python, Git y API. & Construir el flujo y sus componentes. \\
\hline
Pruebas y refinamiento & 24 & pytest, clientes simulados y casos. & Verificar y corregir el sistema y el instrumento. \\
\hline
Evaluación experimental & 20 & OpenAI, Anthropic, SQLite y scripts propios. & Ejecutar tandas y analizar resultados. \\
\hline
Documentación y memoria & 42 & Overleaf, BibLaTeX, repositorio y tutor. & Redactar, revisar y preparar la entrega. \\
\hline
\textbf{Total} & \textbf{180} & \multicolumn{2}{c|}{\textbf{Carga total: 6 ECTS}} \\
\hline
\end{tabular}
\floatfoot{Fuente: elaboración propia.}
\end{table}

La documentación concentra 42 horas porque incluye redacción, bibliografía, tablas, figuras, compilación y revisión. La implementación ocupa 32 horas por abarcar agentes y capas deterministas. Las fases financiera y de pruebas reciben 24 horas cada una debido a la necesidad de contrastar convenciones y revisar tanto el sistema como sus referencias.

\section{Cronograma}
\label{sec:cronograma}

El proyecto comenzó en agosto de 2026 con el análisis y el primer prototipo. El 25 de agosto se alcanzó un producto mínimo viable. Durante septiembre se revisaron el esquema y las convenciones. La reunión con el tutor del 15 de septiembre concretó los calendarios por pata, los siete campos obligatorios, la ubicación de la curva de estimación y el bucle limitado a cinco pasadas.

El 21 de septiembre se cerró el rediseño y se ejecutó la evaluación de seis modelos. El 22 de septiembre concluyó el desarrollo técnico con 84 pruebas y 23 casos verificados. Después se priorizaron la memoria y la defensa para evitar cambios tardíos que invalidaran las tandas archivadas.

\begin{figure}[H]
\centering
\begin{ganttchart}[
    hgrid,
    vgrid,
    x unit=1.55cm,
    y unit chart=0.62cm,
    title height=1,
    bar height=0.55,
    milestone height=0.55,
    title label font=\scriptsize\bfseries,
    bar label font=\scriptsize,
    milestone label font=\scriptsize,
    group label font=\scriptsize\bfseries,
    bar/.append style={fill=gray!45},
    group/.append style={fill=azulUC3M},
    milestone/.append style={fill=black},
    link/.style={->, thick}
]{1}{9}
\gantttitle{Agosto}{4}
\gantttitle{Septiembre}{3}
\gantttitle{Octubre}{2} \\
\gantttitle{S1}{1}\gantttitle{S2}{1}\gantttitle{S3}{1}\gantttitle{S4}{1}
\gantttitle{S1}{1}\gantttitle{S2}{1}\gantttitle{S3}{1}\gantttitle{S1}{1}\gantttitle{S2}{1} \\
\ganttbar{Análisis y alcance}{1}{2} \\
\ganttbar{Prototipo inicial}{2}{4} \\
\ganttmilestone{MVP}{4} \\
\ganttbar{Revisión del esquema}{4}{5} \\
\ganttmilestone{Revisión con el tutor}{5} \\
\ganttbar{Rediseño e implementación}{7}{7} \\
\ganttbar{Evaluación experimental}{7}{7} \\
\ganttmilestone{Cierre técnico}{7} \\
\ganttbar{Redacción de la memoria}{6}{7} \\
\ganttmilestone{Entrega final}{7} \\
\ganttbar{Preparación de la defensa}{8}{9} \\
\ganttmilestone{Presentación}{9}
\end{ganttchart}
\caption[Cronograma del proyecto]{Cronograma de desarrollo, evaluación, redacción y defensa}
\label{fig:cronograma_proyecto}
\floatfoot{Fuente: elaboración propia.}
\end{figure}

El diagrama agrupa las actividades por semanas; por ello, el rediseño y la evaluación realizados el mismo día aparecen en la misma columna. Los hitos técnicos se apoyan en el historial del repositorio y en las tandas archivadas.

\section{Riesgos y decisiones de alcance}
\label{sec:riesgos_alcance}

La Tabla~\ref{tab:riesgos_proyecto} resume los riesgos que condicionaron la planificación y las medidas adoptadas.

\begin{table}[H]
\centering
\small
\caption{Riesgos principales y respuesta adoptada}
\label{tab:riesgos_proyecto}
\begin{tabular}{|P{3.5cm}|P{4.6cm}|P{5.0cm}|}
\hline
\textbf{Riesgo} & \textbf{Impacto} & \textbf{Respuesta} \\
\hline
Ampliación excesiva del dominio & Pérdida de profundidad y casos insuficientes por producto. & Limitar la evaluación a \textit{IRS vanilla} en EUR y USD. \\
\hline
Convenciones incorrectas o desactualizadas & Representación financiera errónea. & Contrastar EUR y USD con fuentes sectoriales y validar lo derivado. \\
\hline
Información ausente o ambigua & Alucinaciones, rechazos o reintentos inútiles. & Separar ausencias no recuperables de errores corregibles; máximo de cinco pasadas. \\
\hline
Defectos del evaluador & Atribuir al modelo cifras falsas pero estables. & Revisar instrucciones, entradas y referencias; conservar respuestas y separar tandas. \\
\hline
Dependencia de proveedores & Cambios de interfaz, parámetros o comportamiento. & Capa común y configuración separada de la lógica financiera. \\
\hline
Cambios tardíos & Invalidación de la evidencia experimental. & Cerrar el código antes de la revisión final de la memoria. \\
\hline
\end{tabular}
\floatfoot{Fuente: elaboración propia.}
\end{table}

La integración con un motor de valoración y la comparación con una arquitectura monolítica se aplazaron para no ampliar el experimento después de cerrar casos y modelos. El prototipo termina en la \textit{RFQ}; no calcula precio ni riesgo. Esta decisión mantiene alineados objetivos, tiempo disponible y evidencia presentada.

La detección de convenciones alternativas mediante un vocabulario cerrado y la imposibilidad de fijar la temperatura se aceptaron como limitaciones conocidas. En ambos casos se priorizó declarar el límite y conservar la trazabilidad frente a introducir cambios de última hora sin una nueva evaluación completa.
